
\documentclass{article}
\usepackage{colortbl}
\usepackage{makecell}
\usepackage{multirow}
\usepackage{supertabular}

\begin{document}

\newcounter{utterance}

\centering \large Interaction Transcript for game `eqbench', experiment `eq\_bench', episode 24 with reflexion{-}Qwen3.5{-}27B.
\vspace{24pt}

{ \footnotesize  \setcounter{utterance}{1}
\setlength{\tabcolsep}{0pt}
\begin{supertabular}{c@{$\;$}|p{.15\linewidth}@{}p{.15\linewidth}p{.15\linewidth}p{.15\linewidth}p{.15\linewidth}p{.15\linewidth}}
    \# & \multicolumn{2}{c}{Player} && \multicolumn{2}{c}{Game Master} \\
    \hline

    \theutterance \stepcounter{utterance}  
    & & & \multicolumn{4}{p{0.6\linewidth}}{
        \cellcolor[rgb]{0.9,0.9,0.9}{
            \makecell[{{p{\linewidth}}}]{
                \texttt{\tiny{[P1$\langle$GM]}}
                \texttt{You have attempted this same task before and did not succeed. These are the plans you wrote after those attempts.} \\
\texttt{Trial \#1: New plan: My previous attempt failed because I assigned a high score to "Angry" (7) despite Eva's dialogue being characterized by cold, professional detachment and intellectual superiority rather than overt emotional outbursts. She explicitly states, "That's my job... I don't coddle," framing her harshness as a professional duty rather than an emotional reaction. To correct this, I will re{-}evaluate the emotional intensity based on her specific tone: I will lower the "Angry" score significantly to reflect her controlled demeanor, while maintaining a high score for "Justified" as she firmly defends her role as a critic. I will also ensure "Apologetic" remains at 0, as she shows no remorse, and keep "Regretful" low, as she appears confident in her actions. The new strategy is to prioritize the emotion that aligns with her stated professional identity over the emotion implied by the conflict's surface tension.} \\
\\ 
\\ 
\\ 
\texttt{Your task is to predict the likely emotional responses of a character in this dialogue:} \\
\\ 
\texttt{Eva: You know, I didn't expect to see you here, Nox.} \\
\texttt{Nox: And I didn't expect you to show up, Eva. Not after everything.} \\
\texttt{Eva: Oh, so you're still playing the victim, huh?} \\
\texttt{Nox: I'm not playing anything, Eva. I'm just trying to move forward.} \\
\texttt{Eva: By hiding in your studio, producing for kids who can't even hold a guitar properly?} \\
\texttt{Nox: At least I'm not tearing them apart in my articles.} \\
\texttt{Eva: That's my job, Nox. I critique. I analyze. I don't coddle.} \\
\texttt{Nox: You always were good at tearing things apart, weren't you?} \\
\texttt{{[}End dialogue{]}} \\
\\ 
\texttt{At the end of this dialogue, Eva would feel...} \\
\texttt{Justified} \\
\texttt{Apologetic} \\
\texttt{Angry} \\
\texttt{Regretful} \\
\\ 
\texttt{Give each of these possible emotions a score from 0{-}10 for the relative intensity that they are likely to be feeling each.} \\
\\ 
\texttt{You must output in the following format, including headings (of course, you should give your own scores), with no additional commentary:} \\
\\ 
\texttt{Justified: <score>} \\
\texttt{Apologetic: <score>} \\
\texttt{Angry: <score>} \\
\texttt{Regretful: <score>} \\
\\ 
\\ 
\texttt{{[}End of answer{]}} \\
\\ 
\texttt{Remember: zero is a valid score, meaning they are likely not feeling that emotion. You must score at least one emotion > 0.} \\
\\ 
\texttt{Your answer:} \\
            }
        }
    }
     \\ \\

    \theutterance \stepcounter{utterance}  
    & \multicolumn{4}{p{0.6\linewidth}}{
        \cellcolor[rgb]{0.561,0.737,0.561}{
            \makecell[{{p{\linewidth}}}]{
                \texttt{\tiny{[P1$\rangle$GM]}}
                \texttt{Justified: 9} \\
\texttt{Apologetic: 0} \\
\texttt{Angry: 2} \\
\texttt{Regretful: 0} \\
            }
        }
    }
    & & \\ \\

    \theutterance \stepcounter{utterance}  
    & & & \multicolumn{2}{p{0.3\linewidth}}{
        \cellcolor[rgb]{0.9,0.9,0.9}{
            \makecell[{{p{\linewidth}}}]{
                \texttt{\tiny{[GM$|$GM]}}
                \texttt{Justified: 9} \\
\texttt{Apologetic: 0} \\
\texttt{Angry: 2} \\
\texttt{Regretful: 0} \\
            }
        }
    }
    & & \\ \\

    \theutterance \stepcounter{utterance}  
    & & & \multicolumn{2}{p{0.3\linewidth}}{
        \cellcolor[rgb]{0.9,0.9,0.9}{
            \makecell[{{p{\linewidth}}}]{
                \texttt{\tiny{[GM$|$GM]}}
                \texttt{Justified: 8} \\
\texttt{Apologetic: 0} \\
\texttt{Angry: 6} \\
\texttt{Regretful: 0} \\
            }
        }
    }
    & & \\ \\

    \theutterance \stepcounter{utterance}  
    & & & \multicolumn{2}{p{0.3\linewidth}}{
        \cellcolor[rgb]{0.9,0.9,0.9}{
            \makecell[{{p{\linewidth}}}]{
                \texttt{\tiny{[GM$|$GM]}}
                \texttt{game\_result = LOSE} \\
            }
        }
    }
    & & \\ \\

    \theutterance \stepcounter{utterance}  
    & & \multicolumn{4}{p{0.6\linewidth}}{
        \cellcolor[rgb]{0.95,0.95,0.95}{
            \makecell[{{p{\linewidth}}}]{
                \texttt{\tiny{[P1$\leftrightarrow$P]}}
                \texttt{Post{-}episode Reflexion (not sent to GM):} \\
\texttt{New plan: My previous attempt failed because I underestimated the intensity of Eva's "Angry" emotion, assigning it a low score of 2 despite her aggressive and confrontational language (e.g., "tearing them apart," "playing the victim"). While she frames her actions as professional ("That's my job"), the tone is clearly hostile and defensive, indicating high underlying anger. I will revise my strategy to recognize that professional detachment in dialogue often masks intense anger in conflict scenarios. Specifically, I will significantly increase the "Angry" score to reflect the aggression in her attacks on Nox's character and work, while maintaining a high "Justified" score for her self{-}defense. I will keep "Apologetic" and "Regretful" at 0, as her dialogue shows no remorse or hesitation. The new plan is to balance the "Justified" score with a much higher "Angry" score to accurately capture the dual nature of her response: defensive justification fueled by active hostility.} \\
            }
        }
    }
    & \\ \\

\end{supertabular}
}

\end{document}
